\uuid{O3Bx}
\titre{Intégration/DL usuels}
\auteur{hanani ; mzari}
\organisation{exo7}
\licence{CC BY-NC-SA 4.0}
\datecreate{2020-01-20}
\niveau{L1}
\module{Analyse}
\chapitre{Développement limité}
\sousChapitre{Calculs}
\difficulte{3}
\theme{Opérations sur les DL, Intégration, DL usuels}
\exo7qcmid{lille-2.developpements-limites.024}

\contenu{
\texte{
Soit $a\in \Rr$ et $f$ une fonction. On dit que $f(x)$ est un $o\left((x-a)^n\right)$ au voisinage de $a$ et on écrit $f(x)=o((x-a)^n)$ si $\displaystyle \lim _{x\to a}\frac{f(x)}{(x-a)^n}=0$. On note par $DL_n(a)f(x)$ le développement limité de $f$ à l'ordre $n$ en $a$. On note aussi $DL_n(a^+)f(x)$ (resp. $DL_n(a^-)f(x)$) le développement limité de $f$ à l'ordre $n$ à droite en $a$ (resp. à gauche en $a$).
}

\begin{question}
\qclassification{125.02}
\qtags{motcle=Intégration/DL usuels}
Soit $\displaystyle f(x)=\arctan (1+x)$ et $\displaystyle g(x)=\arctan (\cos x)$. Parmi les affirmations suivantes, cocher celles qui sont vraies :
\begin{answers}
\good{$\displaystyle DL_2(0)f'(x)=\frac{1}{2}-\frac{x}{2}+\frac{x^2}{4}+o(x^2)$.}
\bad{$\displaystyle DL_3(0)f(x)=\frac{x}{2}-\frac{x^2}{4}+\frac{x^3}{12}+o(x^3)$.}
\good{$\displaystyle DL_3(0)\cos x=1-\frac{x^2}{2}+o(x^3)$ et $\displaystyle DL_3(0)\arctan u=u-\frac{u^3}{3}+o(u^3)$.}
\bad{$\displaystyle DL_3(0)g(x)=\left(1-\frac{x^2}{2}+o(x^3)\right)-\frac{1}{3}\left(1-\frac{x^2}{2}+o(x^3)\right)^3+o(x^3)$.}
\end{answers}
\begin{explanations}
On a : $\displaystyle f'(x)=\frac{1}{1+(1+x)^2}=\frac{1}{2}\times\frac{1}{1+x+\frac{x^2}{2}}=\frac{1}{2}-\frac{x}{2}+\frac{x^2}{4}+o(x^2)$. Ce qui par intégration, puisque $f(0)=\pi/4$, donne
$$DL_3(0)f(x)=\frac{\pi}{4}+\frac{x}{2}-\frac{x^2}{4}+\frac{x^3}{12}+o(x^3).$$
Ensuite, avec $\displaystyle u=\cos x-1=-\frac{x^2}{2}+o(x^3)$, on aura :
$$g(x)=\arctan (1+u)=\frac{\pi}{4}+\frac{u}{2}-\frac{u^2}{4}+o(u^2)=\frac{\pi}{4}-\frac{x^2}{4}+o(x^3).$$
\end{explanations}
\end{question}
}
